\chapter{Literature Review}
\label{ch:lit}

\section{Introduction}

This chapter surveys the research that this project builds upon and
locates the project in the current literature. The survey is
organised into three parts: LLM--based AutoML systems; verification
of generated code and of reasoning; and formal, typed and
graph--based verification of reasoning. The parts are treated
separately because they correspond to different correctness questions
(as formalised in Chapter~\ref{ch:method}) and because each has
developed its own evaluation style. The chapter closes with a
synthesis of the research gap that the proposed framework addresses.

\section{Existing Systems}

\subsection{LLM--Based AutoML Pipelines}

The simplest shape of an LLM--driven AutoML system is a feature
engineering agent that takes a tabular dataset and emits candidate
feature constructions. CAAFE~\cite{caafe} is a representative
instance in which the LLM is used to propose context--aware feature
engineering steps and to iterate on them using data--driven
feedback. LLMFE~\cite{llmfe} extends this idea into an evolutionary
loop in which feature engineering is framed as program search guided
by execution feedback.

Data--centric AutoML is served by CatDB~\cite{catdb}, which uses
data--catalogue--guided prompts to construct ML pipelines, and by
FlowPipe~\cite{flowpipe}, which synthesises data--preparation
pipelines using conditional generative flow networks and
LLM--derived priors over multi--channel preprocessing. Agentic
systems such as DS--Agent~\cite{dsagent} and SELA~\cite{sela} go
further by orchestrating planning, retrieval and search over
candidate pipelines, with case--based reasoning or tree--search
guiding the exploration.

Full--pipeline systems such as AutoML--Agent~\cite{automlagent}
include a multi--agent architecture that spans planning, code
generation and verification. It performs three verification steps:
\emph{request verification} (ReqVer), which assesses whether the
user's instructions are clear and adequate for the ML task;
\emph{execution verification} (ExecVer), which checks --- by
prompting, before any code is run --- whether candidate pipelines
satisfy the user's requirements; and \emph{implementation
verification} (ImpVer), which validates the outcomes of code that has
been compiled and executed. In this way AutoML--Agent explicitly
raises the question of verification inside LLM--driven AutoML, but
these three steps do not establish a one--to--one relationship
between individual plan assertions and program implementations. A
successful implementation--verification verdict therefore does not
entail plan--program parity. The broader survey of LLMs for building
and optimising ML workflows~\cite{llmauto} catalogues this landscape
of systems and the open problems that follow from it.

\subsection{Code and Reasoning Verification}

Program verification has a long history. Classical foundations
include Hoare's axiomatic basis for computer
programming~\cite{hoare1969}, abstract interpretation~\cite{cousot1977}
and model checking~\cite{clarke1999}. These techniques connect
programs to specifications through rules based on assertions,
abstract interpretation, model checking or runtime evidence. They
do not, by themselves, consume LLM reasoning: an LLM plan has to be
lowered into a restricted specification language first, after which
the usual verification machinery applies on representable claims
rather than on universally correct machine learning behaviour.

Executable code benchmarks have shown that behavioural tests give a
stronger signal than surface similarity~\cite{humaneval,
programsynthesis}. However, tests only demonstrate claimed behaviour,
and code generation can leave latent bugs intact~\cite{copilot}. A
survey of LLMs for code generation~\cite{jiang2026survey} catalogues
execution--based benchmarks and lists the integration of formal
verification tools within the LLM pipeline as an open safety
direction, rather than as an existing capability. Empirical studies of
code--generation failure modes~\cite{shen2026failure} offer surface
and runtime error taxonomies alongside decision frameworks that
directly relate failure patterns to actionable strategies in enhancing
performance, including through self--critique, multi--model
collaboration or retrieval--augmented generation. Execution--based
failure, however, is exactly where these frameworks stop; the project
in this report takes over from that point and treats reasoning--code
alignment as the next verification layer.

\subsection{Formal, Typed and Graph--Based Verification of Reasoning}

Chain--of--Thought (CoT) prompting elicits intermediate reasoning
steps~\cite{wei2023cot}. These steps are not required to reflect the
model's underlying computation; recent work has shown that LLMs do
not always say what they think~\cite{turpin2023} and has extended
CoT--faithfulness evaluation from proprietary models to open--weight
reasoning models~\cite{young2026lie}, providing a qualitative basis
for the faithfulness problem. Step--level verifier benchmarks and
process supervision~\cite{jacovi2024chain, lightman2024} inspect
chains at a finer granularity than final answers. These approaches
address the problem of inspectable reasoning mechanisms, which is
known to be difficult for LLMs to resolve on their own~\cite{tyen2024}.

Typed Chain--of--Thought~\cite{perrier2025} synthesises reasoning
steps into typed trees that can be validated, demonstrating that the
tractability of prose can be preserved under constrained
representations, although the resulting certificate is about the
reasoning trace rather than the generated AutoML source code.
CRV~\cite{zhao2026crv} builds an attribution graph of model
computation and predicts the correctness of reasoning based on the
graph structure. In both cases, the verified object is an internal
or near--internal representation of the model, not the executable
pipeline that an AutoML plan is lowered to.

Structured chain--of--thought has been shown to help code
generation~\cite{li2025structured}, but structured generation alone
is not a post--generation proof that all stated steps appear in the
code exactly as stated. Critique and self--refinement are responses
to plan errors or evaluator feedback~\cite{selfrefine, critic}; a
different but executable pipeline can be accepted whenever the
oracle does not return one.

\section{Review of Research Papers}

Table~\ref{tab:relwork} summarises how representative recent work
compares with the proposed framework along five axes: the domain
addressed, the reasoning artifact consumed, the kind of program
artifact produced, whether the system actually executes code,
whether it targets AutoML, and whether it establishes a
mechanical plan--code link.

\begin{table}[H]
\centering
\caption{Representative approaches relative to the proposed
verification layer.}
\label{tab:relwork}
\small
\setlength{\tabcolsep}{4pt}
\renewcommand{\arraystretch}{1.25}
\begin{tabular}{p{2.3cm} p{2.4cm} p{2.4cm} p{2.2cm} c c c}
\toprule
\textbf{Work} & \textbf{Domain} & \textbf{Reasoning artifact} &
\textbf{Program artifact} & \textbf{Exec.} & \textbf{AutoML} &
\textbf{Plan--code link} \\
\midrule
AutoML--Agent~\cite{automlagent} & Full--pipeline AutoML &
Plan guides agents & Implementation run and evaluated & \checkmark &
\checkmark & -- \\

CatDB~\cite{catdb} & Data--centric ML & Catalogue--guided prompts &
Generated pipeline executed & \checkmark & \checkmark & -- \\

FlowPipe~\cite{flowpipe} & Data--prep pipelines & LLM--derived
priors & Policy--generated pipeline & \checkmark & \checkmark & -- \\

Typed CoT~\cite{perrier2025} & Formalised reasoning & Typed
proof/dataflow & Not AutoML code & -- & -- & -- \\

CRV~\cite{zhao2026crv} & Mechanistic verification & Latent
attribution graph & No pipeline code & -- & -- & -- \\

Failure--Aware~\cite{shen2026failure} & Generic code generation &
Absent & Program + error taxonomy & \checkmark & -- & -- \\

Test--based evaluation~\cite{humaneval, programsynthesis} &
Code generation & Usually absent & Program against tests &
\checkmark & -- & -- \\

\textbf{Proposed framework} & LLM--driven AutoML & Atomic typed
claims & AST (pipeline graph proposed) & Complementary & \checkmark &
\checkmark \\
\bottomrule
\end{tabular}
\end{table}

Three patterns emerge. First, systems that execute code either do not
consume plans at all (code benchmarks) or consume them only as
loose context (full--pipeline AutoML agents). Second, systems that
formalise reasoning (typed CoT, CRV) operate on representations that
are not the AutoML program and therefore cannot certify what the
emitted pipeline does. Third, the verification layers that do exist
inside LLM--AutoML systems (for example the three AutoML--Agent
verifiers) check instruction well--formedness, request satisfiability
and runtime success, but not a per--claim correspondence between the
plan and the generated code.

\section{Research Gap}

The synthesis above yields a precise gap: \emph{LLM--driven AutoML
systems can check whether a request is actionable, whether code
runs, and whether outcomes meet thresholds; reasoning--verification
methods can check properties of the rationale; but no reviewed
method provides a code--linked mechanical step that establishes
whether the emitted AutoML code implements the stated plan.}

This matters because, in multi--agent systems, plans are the
artifacts that are presented to users, that record intended
experimental design, that gate human approval and that direct
delegation. A manager can approve one plan while a coding agent
realises another. If the code diverges, the explanation offered to
the user is operationally misleading even when performance is high,
and downstream auditing can be based on a document that no longer
describes the deployed system.

Accordingly, the project in this report sets five research questions
that follow directly from the gap identified above:

\begin{itemize}
    \item \textbf{RQ1:} Can an AutoML plan be transformed into a
    representation mechanically comparable with an executable
    pipeline?
    \item \textbf{RQ2:} Can alignment verification detect
    discrepancies missed by execution--only and performance--only
    checks?
    \item \textbf{RQ3:} Which categories of reasoning--code
    misalignment arise in practice?
    \item \textbf{RQ4:} How does alignment vary with the LLM, the
    reasoning format, the task and the pipeline complexity?
    \item \textbf{RQ5:} Can mismatches be localised before final
    evaluation or deployment?
\end{itemize}

The associated hypotheses are: $H_1$, execution success does not
imply alignment; $H_2$, mechanical alignment discriminates aligned
from misaligned pipelines more reliably than execution success;
$H_3$, mismatches form stable AutoML--specific categories; and
$H_4$, mismatch probability increases with operation count and
dependency depth. These are stated as hypotheses rather than settled
findings. Chapter~\ref{ch:results} reports preliminary evidence from
the 270--cell study: $H_1$ is supported, since seven pipelines
executed and scored well yet failed the faithfulness check. $H_2$ is
illustrated by those same detections but is not quantitatively tested
here --- the head--to--head comparison against execution-- and
performance--only baselines is left to future work, together with the
fuller protocol needed to evaluate $H_3$ and $H_4$ at scale.

\section{Summary}

Existing LLM--AutoML systems demonstrate that end--to--end
pipelines can be built by language models, but their verification
layers focus on request validity, runtime success and headline
performance. Verification of reasoning in the LLM literature
focuses on faithful or structured chains and attribution graphs,
but these are not the AutoML pipeline that is deployed. Classical
program verification offers the formal apparatus for
specification--based reasoning, but requires a specification --- which
an LLM plan does not readily provide. The combination of these
observations motivates a dedicated \emph{reasoning--code alignment}
verification layer for LLM--driven AutoML, which is the subject of
the next chapter.
